\section{Introduction}

Reading-time experiments are intended to distinguish accounts of how people
understand language. The observations used to make that distinction are,
however, not a direct record of a single cognitive process. A long response
can reflect difficult integration, an unsuccessful retrieval, a momentary
interruption, or a delay unrelated to the sentence. Very short responses can
reflect efficient recognition, anticipation, or an accidental keypress.
Researchers must decide which observations to retain, on what scale to model
them, and how to represent variation between readers and materials. These
decisions are necessary, but their consequences are not necessarily negligible.

Outlier treatment and transformation address different problems. Removing a
response changes the set of events represented in the analysis; transforming
a response changes the scale on which retained events are compared. Neither
operation is a substitute for the other. A fixed upper bound, for example,
removes the same physical duration in every condition, whereas a
distribution-relative bound can remove a response in one condition and retain
an equally long response in another. A logarithm compresses the upper tail
without deleting it. Consequently, agreement between two analyses that both
use ``trimmed RTs'' cannot be assumed merely from that description.

There are substantive reasons to examine these decisions in sentence
processing. In self-paced reading, readers advance through successive words
or phrases, and comprehension is usually assessed separately. In a Maze
task, each advance requires choosing the next sentence segment against a
distractor. This additional decision can change both the distribution of
durations and which observations become available after an error
\citep{boyce_maze_2020}. Even within a task, different constructions and
regions need not have interchangeable baselines. A threshold appropriate
for one region may remove a meaningful part of another region's response
distribution. The general problem predates current web-based experiments:
\citet{ratcliff_methods_1993} demonstrated why the handling of response-time
distributions deserves methodological attention.

The relevant question is therefore not simply whether a particular pipeline
produces a significant effect. It is whether conclusions, estimated effects,
and uncertainty remain similar over an explicitly defined set of defensible
pipelines. An analysis that reports only the preferred pipeline does not show
this dependence. Researcher flexibility can affect reported evidence
\citep{simmons_false-positive_2011}; a multiverse makes a specified part of
that flexibility visible rather than allowing it to remain implicit
\citep{steegen_multiverse_2016,simonsohn_specification_2020}.

Existing methodological evidence motivates, but does not replace, a
sentence-processing reanalysis. Simulation studies demonstrate that the
performance of RT exclusion methods depends on the assumed contamination
process \citep{ratcliff_methods_1993,berger_outliers_2021}. More recently,
\citet{bognar_multiverse_2025} used a multiverse simulation of the congruency
sequence effect to examine how analytical pathways affect Type I and Type II
errors. The present study asks a complementary question: how much do
inference, estimated effects, and precision actually vary across archived
SPR and Maze inputs under a common declared workflow? Unlike a simulation,
this empirical corpus cannot establish which pipeline recovers the truth.

The present analysis crosses 13 exclusion rules with six transformation
specifications in a source-audited collection of existing reading experiments.
The initial source bibliography contains approximately 80 experiment entries.
The implemented analysis is a source-audited collection of analytical inputs,
not a claim to have reanalysed 80 independent experiments. We distinguish inputs for which
untrimmed RTs could be reconstructed or identified from author-prepared
inputs for which earlier exclusions cannot be reversed. This distinction
matters: a permissive new cutoff cannot restore observations that an earlier
cutoff already removed.

We ask three related questions. First, how often does a coefficient's
classification relative to \(p=.05\) change across preprocessing pipelines?
Second, how much do effect estimates, their directions, and standard errors
vary? Third, do the observed patterns appear to differ with the task,
initial distribution, or dataset size? These are questions about analytical
dependence in the supplied data, not tests of which pipeline recovers an
unknown true effect. The multiverse does not turn a majority of significant
pipelines into independent replications or into a probability that a
linguistic hypothesis is correct.

\section{Preprocessing procedures}

\subsection{Common eligibility and comparison subsets}

Before constructing the multiverse, we identify the response column,
participant and material identifiers, factorial manipulation, and analysis
regions from the supplied files and analysis scripts. Missing, nonfinite,
and nonpositive RTs cannot be analysed by every transformation and are
excluded from every pipeline. Practice, filler, and questionnaire records
are excluded where the adapter identifies them. Documented non-RT
eligibility decisions are held constant where they can be reconstructed.
The Saul DO adapter also holds fixed the authors' two participant-level
exclusions for RT-based reasons; ``untrimmed'' in that case describes the
available trial RTs conditional on that participant restriction.
These initial decisions are recorded separately from the observations
subsequently removed by a multiverse rule.

Critical sites are identified from source reports and scripts. A post-critical
site is included only where the source paper reports an experimental main
effect there; examining a site, finding an interaction alone, or finding a
nuisance age effect does not suffice. Distinct critical sites and documented
post-critical positions remain separate analysis jobs. Region selection is
therefore source-based, not selected from the new multiverse results. For the
Boyce datasets, the authors' supplied relative-clause label join identifies
the disambiguating position with its original SPR offset. For ellipsis,
the source-defined \emph{did too} unit is summed rather than treating
\emph{too} alone as the complete region. Other multi-position definitions
follow the source-specific adapter and are documented in the region ledger;
no adapter chooses a region by its new significance or its fastest response.

All distribution-relative rules are calculated within the same
condition--region subset, pooling observations across participants and
items. The focal participant's observations remain in that reference set.
The reference statistic is calculated once from the eligible input; removal
is not iterated. This common grouping isolates the choice of reference
statistic and bounds. It is deliberately not a claim that every original
article used this grouping. In particular, a participant-relative SD rule
from an original analysis is not reproduced merely by including an SD rule
in the present multiverse.

\subsection{Exclusion rules}

The 13 rules comprise no additional exclusion; mean \(\pm k\) sample SD for
\(k\in\{1.5,2,2.5,3\}\); fixed limits of 100--2,000, 100--3,000,
200--2,000, or 200--3,000 ms; the median-relative rule implemented in
\rtprocess; and median \(\pm k\) MAD for \(k\in\{2,2.5,3\}\).
Boundary observations are retained: exclusion requires a value strictly
outside a bound. ``No exclusion'' means no additional RT-based removal after
common eligibility, not the recovery of missing or previously deleted data.

For the median-relative rule, an observation \(x\) is retained when
\[
            0.5M \leq x \leq 2.8M,
\]
where \(M\) is the subset median. This is equivalently
\(M-0.5M \leq x \leq M+1.8M\). The distinction between an additive
distance of \(1.8M\) above the median and an upper multiplier of \(1.8\)
is important: the implemented upper bound is \(2.8M\).
The MAD rules use the normal-consistency-scaled median absolute deviation,
\[
       \operatorname{MAD}=1.4826\,
          \operatorname{median}_i|x_i-M|.
\]
Thus, the numerical MAD multipliers do not refer to an unscaled absolute
deviation. Explicitly stating that constant is necessary for reproducing
the comparison.

The rules represent different assumptions, not different strengths of a
single intervention. Fixed bounds assume that physical duration is the
relevant criterion. SD rules depend on the mean and spread, both of which
can be influenced by an upper tail. Median-relative bounds tie the
permitted interval to a robust measure of location but not to an
independently estimated spread. MAD rules use robust location and spread.
No rule, by itself, establishes that every excluded event is cognitively
impossible. We therefore report removal proportions and resulting inference
rather than defining success as the greatest amount of removal.

\subsection{Transformations}

After exclusion, each retained dataset is represented on six requested
scales: milliseconds, natural logarithm, base-10 logarithm,
\(-1000/x\), square root, and Box--Cox
\citep{box_analysis_1964}. The inverse is increasing over positive RTs,
as are the other transformations used here, so the sign of a simple
condition difference is not mechanically reversed by defining
\(-1000/x\) rather than \(1/x\). Nevertheless, differences between
distributional locations need not remain equal across nonlinear scales.

For Box--Cox,
\[
 g_\lambda(x)=
 \begin{cases}
  (x^\lambda-1)/\lambda,&\lambda\neq0,\\
  \log(x),&\lambda=0.
 \end{cases}
\]
We estimate one \(\lambda\) for each dataset--region--exclusion combination,
not a separate transformation for each condition. Estimation minimizes the
Gaussian profile objective from the complete fixed-effects design, including
the Jacobian, over \([-2,2]\). This is a fixed-design screening profile,
not a full mixed-model likelihood profile. Boundary estimates are recorded
and should not be mistaken for unconstrained optima. The purpose of varying
\(\lambda\) is to allow the shape adjustment to depend on the retained data,
without allowing conditions within the same comparison to have incompatible
measurement scales.

Natural and base-10 logarithms satisfy
\(\log_{10}(x)=\log(x)/\log(10)\). They consequently produce identical
tests in otherwise identical Gaussian mixed models; their coefficients and
standard errors differ only by a known constant. We include both in the
78 requested specifications but obtain the base-10 output by exact
rescaling of the natural-log fit. Summaries of analytical disagreement use
the 65 distinct exclusion--scale combinations, not 78 purportedly independent
choices.

\section{Source datasets}

\input{generated/corpus_text.tex}

Primary inputs retain the original RT column or are reconstructed before the
author's trial-level RT trimming, conditional on disclosed common eligibility
restrictions. Previously prepared secondary inputs remain in the audit ledger
but are not pooled with the primary run or presented as a completed secondary
analysis. Their distributions may already reflect unrecoverable decisions.
This distinction is about provenance, not the reliability of the source
research. The inventory reports paper-derived preprocessing policies rather
than substituting a script filename for a description. Discrepancies between
paper and script are preserved in the accompanying policy evidence table.

The collection is heterogeneous and is not a random sample of
psycholinguistic experiments. Multiple datasets can come from the same
publication; split manipulations may reuse participants; and critical and
post-critical regions belong to the same experimental observations.
Summed participant counts would therefore overstate the number of
independent people. Dataset-specific counts are provided instead. Claims
about task differences are necessarily limited by the number and identity
of source publications, not just by the number of fitted models.

\section{Analysis}

\input{generated/workflow.tex}
Figure~\ref{fig:workflow} summarizes the full workflow. A specification is
the combination of screening rule and outcome scale, together with the
fixed model-reduction algorithm; it is not just a transformation applied
to a single unchanging fitted model.

\subsection{Inferential models and random-effects reduction}

For each dataset and region, the fixed-effects design contains the
documented experimental factors and their interactions. Binary factors are
sum-coded \(-0.5,+0.5\); factors with more levels use \(K-1\) explicit
sum contrasts scaled by one half. The latter are not interpreted as
pairwise tests merely because individual coefficients are printed. We
retain the third factor in the Tucker experiment that also manipulates
ambiguity, rather than collapsing it into an unexplained condition label.
Factor-level mappings accompany the reproducibility outputs.

We fit Gaussian linear mixed models with \texttt{lme4}
\citep{bates_fitting_2015}, beginning with random intercepts for readers
and materials and all design-supported slopes
\citep{barr_random_2013}. A slope is included only when all factors
participating in that slope vary within the corresponding grouping unit.
Thus, the instruction to fit a maximal model does not require an
unidentifiable within-participant slope for a between-participant factor.
Numeric contrast columns make the independent-slope specification explicit,
including for interactions and factors with more than two levels.

If the initial model fails the acceptance criteria, we first remove random
effect correlations. We then remove the random component with the smallest
estimated variance and refit, repeating until the fit is both converged
and nonsingular. The procedure considers intercepts as well as slopes;
any removal of a grouping intercept is recorded rather than hidden.
If no estimable random structure remains, or a fit error prevents an
estimated-variance ordering, the pipeline is reported as unsuccessful.
We do not declare success by replacing a failed mixed model with an
independent-observations regression.

Models use REML and the \texttt{nloptwrap} optimizer, with at most
20,000 objective evaluations per attempt. Source-verified unchanged
critical-region caches from an earlier run retain their original
\texttt{bobyqa} fits, with the engine identified in each checkpoint.
Corrected inputs cannot reuse those caches. Singularity is assessed
at tolerance \(10^{-4}\). Preliminary singular fits need not undergo
expensive numerical derivative calculations; every provisionally accepted
fit is refitted with derivatives enabled and must have optimizer code zero,
no \texttt{lme4} convergence message, and a nonsingular random-effects
covariance estimate. The audit retains formulas, warnings, components
removed, estimated removal variances, and final variances.

For numerical stability, the transformed response is centered and divided
by its within-pipeline SD during fitting. Reported native-scale coefficients
and SEs are restored algebraically. We obtain approximate Satterthwaite
degrees of freedom and two-sided \(t\)-distribution \(p\) values using
\texttt{lmerTest}; 95\% intervals use the corresponding \(t\) quantile.
An unadjusted \(p<.05\) defines the classification whose stability is
examined. It does not designate one pipeline as the preferred analysis.
Significant interactions receive conditional simple-effect contrasts of
the accepted model, with its covariance and degrees-of-freedom calculation
retained. These exploratory follow-ups are reported separately from the
unconditionally evaluated fixed coefficients.

\subsection{Stability outcomes}

We distinguish significance, effect, and precision stability. A coefficient
is significance-unstable if at least one successful distinct pipeline gives
\(p<.05\) and another does not. We also report the proportion of successful
pipelines that disagree with a fixed reference: no additional exclusion
followed by a natural logarithm. That reference supplies a common
comparison, not ground truth or a reconstruction of every source article.
Failed models are shown separately and are not silently treated as
nonsignificant results.

For effect stability, we retain native-scale estimates, direction,
and standardized estimates obtained by dividing a coefficient by the SD of
its pipeline's transformed outcome. The standardized quantity permits a
dimensionless description, but its denominator also changes with
preprocessing. Accordingly, a change in it is not automatically a change in
one invariant psychological effect. We do not directly pool an
inverse-scale coefficient with a coefficient in milliseconds. Precision
is represented by native and standardized SEs and by 95\% intervals.
Direction changes and movement across \(p=.05\) are separate outcomes:
a coefficient can remain positive while changing significance.

Random-effects reduction can itself affect precision. We therefore record
whether each accepted formula matches the reference formula and provide
disagreement summaries restricted to that same-structure subset. These
summaries help identify dependence that is not solely a consequence of a
different final formula; they do not make model adaptation ignorable.
The principal estimand is the result of a complete specified workflow,
including the required convergence procedure.

We additionally visualize changes in estimates and uncertainty for every
exclusion--transformation combination relative to no additional exclusion
and untransformed RTs. This separate reference makes exclusion alone,
transformation alone, and their combination directly identifiable; it does
not replace the natural-log reference for significance stability.
Fixed-effect shifts are absolute differences between standardized coefficients.
SE and confidence-interval-width changes are absolute percentage differences
from their corresponding standardized reference values, so increases and
decreases do not cancel. Interaction-magnitude shifts compare the absolute
values of corresponding standardized interaction coefficients.

Random-effect estimate shifts average absolute differences in shared
variance and covariance entries, expressed in pipeline outcome-variance
units. Entries are matched by grouping factor and named design terms, not
by their printed row order. Removed or added entries are not assigned an
estimated zero. Instead, structural change is shown separately as
\(100[1-|A\cap B|/|A\cup B|]\), where \(A\) and \(B\) are the sets of
estimated covariance entries in the two accepted models. These summaries
use the same publication weighting as the other descriptive outcomes.

\subsection{Isolating exclusion and transformation decisions}

A workflow that differs from the reference in both exclusion and scale
does not reveal which decision produced the difference. We therefore
also make matched comparisons in which only one preprocessing dimension
changes. For exclusion, we compare each of the 78 distinct pairs of rules
on each of the five nonredundant scales. For transformation, we compare
each of the ten distinct pairs of scales under each of the 13 exclusion
rules. The same coefficient and region appear on both sides of every
comparison. A changed classification means that exactly one of the two
tests has \(p<.05\). We also calculate the absolute difference between
their signed \(t\) statistics.

We average these changes across the levels of the dimension held fixed,
then across the distinct pairs of the dimension being changed. An
individual rule's sensitivity is its mean disagreement with each other
rule, averaging over scales; an individual scale's sensitivity is its
mean disagreement with each other scale, averaging over rules. These
summaries do not require choosing one preferred workflow as an anchor.
We retain anchored comparisons with no exclusion or natural logarithms
for ease of interpretation. Dataset and publication weighting follow
the procedure below. For the exclusion-minus-transformation comparison,
the bootstrap resamples the same publications for both quantities.

These are descriptive comparisons over a finite, deliberately enumerated
decision space, not percentages of variance explained or causal
attributions. The numbers of levels and their spacing differ across
dimensions. Each workflow retains its own accepted random-effects
structure, so changing only preprocessing does not require an unchanged
final covariance model. Matched-pair summaries also record whether the
accepted formulas agree.

\subsection{Aggregation and exploratory moderators}

Pipeline outputs share nearly all their observations, so they are not
independent replications. We first summarize within a coefficient and
region, then within a dataset, and finally give source publications equal
weight. Uncertainty intervals for those summaries resample entire
publications 1,000 times. This preserves the grouping of split datasets and
regions within a source, although it cannot supply missing source-level
diversity.

Associations with initial RT skewness, participant count, item count, and
absolute reference-effect magnitude are
exploratory Spearman correlations between dataset-level disagreement
summaries and the corresponding characteristics. Their uncertainty is
estimated by the same publication-cluster resampling. Task comparisons are
descriptive: a task is not randomly assigned across independent source
publications in this corpus. No causal interpretation or universal ranking
of preprocessing policies follows from these associations.

\section{Results}
\input{generated/results_text.tex}

\begin{figure*}[t]
\centering
\includegraphics[width=.93\textwidth]{figures/retention.pdf}
\caption{Additional exclusion under each rule in primary inputs. Points average
regions within analytical inputs and then give source publications equal weight;
intervals resample publications 1,000 times. Fixed-bound labels are in ms.
The median rule retains \(0.5M\)--\(2.8M\); MAD is consistency-scaled.
The figure describes the supplied region-level eligible observations,
not the proportion originally excluded by source authors.}
\label{fig:retention}
\end{figure*}

\begin{figure*}[t]
\centering
\includegraphics[width=.93\textwidth]{figures/inference_stability.pdf}
\caption{Source-weighted percentage of successful coefficient--region tests
whose \(p<.05\) classification differs from the no-additional-exclusion,
natural-log reference in primary inputs.
Each cell gives the percentage numerically as well as by shading.
Natural and base-10 logs are identical for this outcome, so only one is shown.
The surface describes complete workflows; Figure~\ref{fig:decisions}
separates changes in exclusion from changes in transformation.
Cells are not independent evidence.}
\label{fig:inference}
\end{figure*}

\begin{figure*}[t]
\centering
\includegraphics[width=\textwidth]{figures/decision_sensitivity.pdf}
\caption{Inferential sensitivity when one decision dimension changes at a time.
A: percentage of matched comparisons whose \(p<.05\) classification changes
when exclusion changes at fixed scale, or scale changes at fixed exclusion.
B: each rule's average disagreement with the other 12 rules, across five
scales. C: each scale's average disagreement with the other four scales,
across 13 rules. Rules and scales are ranked by these descriptive rates;
the panels do not show statistical accuracy or variance explained.
All panels give publications equal weight after within-input averaging;
intervals are 95\% publication-cluster bootstrap intervals. The same
horizontal scale is used throughout. Fixed bounds are in milliseconds,
the median rule retains \(0.5M\)--\(2.8M\), and MAD is consistency-scaled.}
\label{fig:decisions}
\end{figure*}

\begin{figure*}[t]
\centering
\includegraphics[width=.93\textwidth]{figures/lambda.pdf}
\caption{Distribution of estimated Box--Cox parameters in primary
dataset--region combinations after each exclusion rule. The dashed line
marks the logarithmic limit. Estimates come from a fixed-design profile
over \([-2,2]\); estimates near either endpoint are flagged in the
reproducibility outputs. Boxes describe between-input variation, not
confidence intervals for a common parameter.}
\label{fig:lambda}
\end{figure*}

\begin{figure*}[t]
\centering
\includegraphics[width=.94\textwidth]{figures/precision.pdf}
\caption{Effect and precision stability across distinct successful workflows.
Left: within-coefficient range of standardized estimates. Right:
base-2 logarithm of the ratio of largest to smallest standardized SE; 0
denotes equal SEs, 1 a twofold ratio, and 2 a fourfold ratio. Each
coefficient--region combination contributes once. These are descriptive
distributions, not confidence intervals or pipeline-independent replications.}
\label{fig:precision}
\end{figure*}

\clearpage
\onecolumn
\unsetvruler
\begin{landscape}
\begin{figure}[p]
\centering
\includegraphics[width=.94\linewidth]{figures/preprocessing_effects.pdf}
\caption{Changes in estimates and uncertainty across all requested preprocessing
combinations, relative to no removal and untransformed RTs. Outlined rows and
columns identify transformations alone and exclusions alone. A--E summarize
absolute changes; F distinguishes structural changes from shifts in shared
random-effect estimates. Standardization removes arbitrary units, not
differences between nonlinear-scale estimands. Publication-bootstrap intervals
and signed SE/interval-width changes accompany the reproducibility materials.}
\label{fig:estimate-sensitivity}
\end{figure}
\end{landscape}
\clearpage
\twocolumn
\linenoon

\section{Discussion}

\input{generated/discussion_text.tex}

\subsection{Interpreting analytical dependence}

The contribution is not a claim that preprocessing flexibility is newly
discovered. Rather, it is a source-audited empirical characterization of
that flexibility in sentence-level reading tasks. In multiverse analysis,
alternative datasets and workflows make analytical dependence visible
\citep{steegen_multiverse_2016}; specification curve analysis additionally
emphasizes theoretically justified, statistically valid and nonredundant
specifications \citep{simonsohn_specification_2020}. Counting two logarithmic
bases as separate inferential choices would violate the nonredundancy
principle. Even after collapsing them, however, our uniformly enumerated
65 workflows should not be interpreted as equally plausible scientific
models, and their fractions are not posterior probabilities.

This distinction also separates the present study from the simulation
framework of \citet{bognar_multiverse_2025}. Their known generating
conditions permit examination of false-positive and true-positive rates.
Here, the true status and size of source effects are unknown. Agreement
with the no-exclusion logarithmic reference is therefore an index of
analytical similarity, not estimated accuracy, power, replication success,
or Type I error control. The reported confidence intervals describe
between-publication uncertainty in this corpus; they are not joint
specification-curve tests of a common linguistic hypothesis.

The condition-relative rules require particular care. A distributional
shift caused by a manipulation can affect a condition's own reference
median, mean, or spread. This may be appropriate when the goal is to remove
responses that are unusual relative to a local baseline. It is not
equivalent to enforcing a common duration threshold across conditions.
Our common condition--region grouping makes that choice explicit but
does not establish that it is always preferable to participant-relative
or model-residual screening. Those grouping decisions are additional
dimensions of a larger multiverse.

\subsection{Exclusion, transformation, and the scientific estimand}

The present results should be read alongside, rather than as a replacement
for, research on identifying invalid responses. \citet{berger_outliers_2021}
showed that exclusion methods behave differently when simulated contaminants
occupy the tails versus overlap the valid distribution. Thus, no empirical
ranking by retention or significance is sufficient to identify a
contamination-removal rule. \citet{miller_pooling_2024} developed an empirical
CDF-pooling approach to characterize outliers in large lexical-decision
datasets. That work illustrates the value of independent information about
contaminants; its estimates cannot simply be transferred to contextual
reading or Maze decisions. Our removed observations were not labelled as
cognitively invalid, and the study supplies no such labels retrospectively.

Transformation raises a different scientific issue. \citet{lo_transform_2015}
demonstrated that nonlinear transformations can alter conclusions about
additivity and interaction, and discussed generalized mixed models as an
alternative that makes distributional and link assumptions explicit.
Consequently, a significant raw-scale interaction and a nonsignificant
log-scale interaction need not be competing answers to exactly the same
estimand. Our standardized ranges facilitate numerical description but
do not eliminate this distinction. A theory of absolute processing
delays, proportional slowing, or changes confined to the upper tail may
require different models and different summaries.

The variation in Box--Cox parameters therefore motivates explicit scale
assessment, not automatic preference for whichever scale maximizes
normality or significance. The profile used here concerns the fixed-effects
Gaussian screening model and is not a substitute for checking the final
model's residual structure, heteroscedasticity, and dependency assumptions.
Generalized, distributional, robust or mixture models are reasonable
extensions when motivated by the task and estimand; they were not fitted
in this multiverse and are not claimed to outperform its workflows.
As \citet{vankov_hazards_2023} emphasizes, switching to a more flexible model
does not by itself remove researcher degrees of freedom.

\subsection{Scope and limitations}

Several limitations qualify generalization. First, the analysis covers an
audited convenience corpus rather than a probability sample of the literature.
One computationally demanding primary input was excluded; this runtime-based
exclusion could underrepresent difficult fitting problems. Unanalysed inputs
must not be read as experiments with null effects. Second, earlier eligibility
decisions can affect the available distribution even when original RT columns
survive. Previously prepared inputs are inventoried but excluded from this
primary reanalysis, and known participant-level restrictions are disclosed. Third,
correct-response Maze data are conditional on an earlier sequence of
successful decisions. No RT transformation restores regions that were
never presented or responses that were not saved.

Fourth, the fixed-design Box--Cox profile is a practical screening
approximation, not optimization of the final mixed-model likelihood.
Fifth, variance-guided random-effects reduction changes the fitted
covariance model when preprocessing changes what the design supports.
Its trace is part of the result. A converged and nonsingular fit is a
necessary operational acceptance criterion here, not proof that the
model is a complete account of trial-level dependence or that every
Gaussian assumption holds.

Finally, the multiverse is finite and investigator-defined. It does not
cover every possible cutoff, winsorization procedure, residual screen,
distributional model, or grouping choice. In particular, the two
logarithms cannot be counted as independent specifications merely to
increase the apparent size of the analysis. A different set of plausible
pipelines could yield different numerical proportions. Our aim is a
reproducible description of the declared decision space, not an estimate of
the probability of every analysis that a researcher might have performed.

\subsection{Practical advice and recommendations}

The results support transparent specification rather than a universal
winning cutoff. We recommend the following practices for new reading-time
studies and secondary analyses.

\begin{enumerate}
\item \textbf{Define the scientific target before choosing the scale.}
State whether the hypothesis concerns absolute delay, proportional slowing,
a factorial interaction on a specified scale, or a distributional feature.
Choose a primary analysis on that basis, not on the smallest observed
\(p\) value. Nonlinear alternatives should be described as
scale-dependent sensitivity analyses, not exact replications of one estimand.

\item \textbf{Specify eligibility separately from outlier screening.}
Document task errors, post-error handling, missing responses, participant
restrictions, and source-defined analysis regions. For each RT rule give
the physical or relative bounds, comparison subset, MAD constant where
applicable, boundary convention, iteration policy, and order relative
to transformation. Preserve original RTs and an auditable exclusion flag.

\item \textbf{Declare a limited, justified sensitivity set.}
Preregister the primary workflow and a scientifically defensible set of
alternatives where possible. Remove algebraically redundant specifications.
In this corpus, exclusion rules, especially the 1.5-SD rule, and the inverse
scale produced the largest inferential changes in their respective
dimensions. They deserve explicit attention when scientifically plausible;
comparing natural and base-10 logarithms adds no inferential information.
The relatively low average disagreement between logarithms and Box--Cox
does not guarantee equivalence for a particular interaction.
Do not choose a workflow after seeing which supports the intended account:
simulation evidence shows that such flexibility can affect significance,
confidence intervals, and effect-size reporting \citep{vankov_hazards_2023}.
A transparent multiverse documents uncertainty; it is not a license to
selectively report the preferred branch.

\item \textbf{Report three forms of stability.}
Show native-scale estimates, SEs and intervals for the primary model;
summarize direction and effect magnitudes across alternatives; and distinguish
threshold crossings from substantial changes in interpretation. A fraction
of significant workflows is not a substitute for a joint test, a replication,
or a probability that the effect exists.

\item \textbf{Make model adaptation and reproducibility visible.}
State the maximal supported random structure and reduction rule, save
convergence and singularity diagnostics, and show any removed intercepts.
A nonsingular model is not automatically scientifically adequate. Supply
compact original-scale inputs, factor mappings, source and region evidence,
full pipeline outputs, and scripts that regenerate the manuscript numbers
and figures. Keep already prepared inputs distinct from recoverable
pre-exclusion RTs, and disclose computational exclusions.
\end{enumerate}

Together, these practices make the substantive claim, the measurement
scale, and the effect of analytical choices inspectable. That is the
practical benefit of this multiverse: not selecting a generally optimal
preprocessing rule, but showing which conclusions require a more
qualified account of how the data were analysed.

\section{Reproducibility and declarations}

The source-specific adapters, multiverse runner, model-reduction audit,
coefficient outputs, conditional contrasts, and figure generators accompany
the manuscript. Original source files are not overwritten. Original-data
hashes identify the exact local inputs, and all regenerated outputs are
written separately. The full run is deterministic apart from numerical
optimization; the publication-cluster bootstrap uses a fixed random seed.
Compact primary RDS inputs permit model reproduction without searching
archives again. Completed checkpoints permit immediate regeneration of tables
and figures without refitting. \texttt{README.md} lists the extraction,
fitting, verification, and manuscript-build commands. Runtime depends strongly
on input and covariance complexity; some region jobs required many hours%
\input{generated/runtime_footnote.tex}.

This study reanalyses existing data and recruited no new participants.
Source citations and host identifiers are reported in the dataset tables.
Funding and competing-interest declarations require author confirmation
before submission; no declaration is inferred from the reanalysis itself.

\clearpage
\onecolumn
\unsetvruler
\begin{landscape}
\section*{Dataset inventory}
\input{generated/dataset_tables.tex}
\end{landscape}
\clearpage
\twocolumn
\linenoon
\bibliography{multiverse_refs}
